% ตัวอย่างวิทยานิพนธ์ภาษาไทย -- compile:  latexmk -xelatex example-th
\documentclass[thai]{kuthesis}
% options: thai | english, numbered (เลขกำกับหัวข้อข้าง), guides (เส้นวัดระยะ),
%          latin=inter|sarabun, math=fira|lete|newcm|stix|none, thaifont=...

\usepackage[style=kuthesis]{biblatex}
\addbibresource{references.bib}

\thesistitle{การบำบัดน้ำเสียด้วยจุลินทรีย์แบบยึดเกาะตัวกลางชนิดเส้นใยไนล่อนโดยระบบถังกรองไร้อากาศ}
            {Wastewater Treatment of Fixed Film Microorganism on Nylon Rope Media by Anaerobic Filter}
\thaiauthor[นางสาว]{ซันวานี}{จิใจ}
\englishauthor[Miss]{Sunwanee}{Jijai}
\degree{วิศวกรรมศาสตรมหาบัณฑิต (วิศวกรรมสิ่งแวดล้อม)}
       {Master of Engineering (Environmental Engineering)}
\majorfield{วิศวกรรมสิ่งแวดล้อม}{Environmental Engineering}
\department{วิศวกรรมสิ่งแวดล้อม}{Environmental Engineering}
\advisor{รองศาสตราจารย์สมชาย ใจดี, Ph.D.}{Associate Professor Somchai Jaidee, Ph.D.}
\coadvisor{ผู้ช่วยศาสตราจารย์สมหญิง รักเรียน, D.Eng.}{Assistant Professor Somying Rakrian, D.Eng.}
\departmenthead{รองศาสตราจารย์สมศักดิ์ ตั้งใจ, Ph.D.}{Associate Professor Somsak Tangjai, Ph.D.}
\dean{รองศาสตราจารย์กัญจนา ธีระกุล, D.Agr.}{Associate Professor Gunjana Theeragool, D.Agr.}
\thesisyear{2549}{2006}
\ackdate{สิงหาคม 2549}{August 2006}
\kulogo[logo/ku-logo-en.png]{logo/ku-logo-th.png}  % ตรามหาวิทยาลัย (ไม่มีไฟล์ = วงกลมแทน)
\kuapprovallogo{logo/ku-logo-th-green.png}

\begin{document}

\frontmatter
\makecover
\makeapproval
\maketitlepage

\begin{thaiabstract}
  งานวิจัยนี้ศึกษาการบำบัดน้ำเสียด้วยถังกรองไร้อากาศ (anaerobic filter) ที่ใช้เส้นใยไนล่อนเป็นตัวกลางให้จุลินทรีย์ยึดเกาะ
  โดยเดินระบบที่ภาระบรรทุกสารอินทรีย์ 0.5--4.0 kg COD/m\textsuperscript{3}$\cdot$d
  ผลการทดลองพบว่าระบบสามารถกำจัดซีโอดีได้ร้อยละ 85--92

  ระยะเวลากักเก็บน้ำที่เหมาะสมคือ 12 ชั่วโมง และอัตราการเกิดก๊าซชีวภาพแปรผันตามภาระบรรทุกสารอินทรีย์
\end{thaiabstract}

\begin{englishabstract}
  This research studied an anaerobic filter using nylon rope as the support medium for
  attached microorganisms. The reactor was operated at organic loading rates of
  0.5--4.0~kg~COD/m\textsuperscript{3}$\cdot$d and removed 85--92\% of COD.

  The optimum hydraulic retention time was 12 hours.
\end{englishabstract}

\begin{acknowledgments}
  ผู้วิจัยขอกราบขอบพระคุณ รองศาสตราจารย์ ดร.สมชาย ใจดี ประธานกรรมการที่ปรึกษา
  ที่ให้คำปรึกษาในการเรียน การค้นคว้าวิจัย และการเขียนวิทยานิพนธ์จนเสร็จสมบูรณ์

  ขอขอบพระคุณบิดา มารดา ที่ให้การอบรมและให้กำลังใจผู้วิจัยมาตลอด
\end{acknowledgments}

\tableofcontents
\listoftables
\listoffigures

\begin{abbreviations}
  \abbr{BOD}{biochemical oxygen demand}
  \abbr{COD}{chemical oxygen demand (ความต้องการออกซิเจนทางเคมี)}
  \abbr{HRT}{hydraulic retention time}
  \abbr{OLR}{organic loading rate}
  \abbr{$\mu$}{อัตราการเจริญจำเพาะ (d\textsuperscript{-1})}
\end{abbreviations}

\mainmatter

\chapter{คำนำ}

ผีเสื้อจัดเป็นแมลงในอันดับ (Order) Lepidoptera แบ่งได้เป็น 2 ประเภท คือ ผีเสื้อกลางวัน
และผีเสื้อกลางคืน \textcite{boddington2000a} กล่าวว่าการจัดการทางการเกษตรมีผลต่อเชื้อรา
อาร์บัสคูลาร์ไมคอร์ไรซา ขณะที่ \textcite{patcharaporn2552} เสนอตัวแบบการจัดสรรทรัพยากร
และมีการศึกษาที่เกี่ยวข้องอีกหลายเรื่อง \parencite{siriyaporn2552,chutima2551,peterson2004}

การอ้างอิงเฉพาะหน้าใช้เครื่องหมายทวิภาค เช่น \parencite[2-5]{pitchapa2552}
หรือ \textcite[3]{bago2000} ส่วนข้อมูลจากเว็บไซต์อ้างอิงเช่นเดียวกัน \parencite{moac2551,who2006}

\chapter{วัตถุประสงค์}

\begin{enumerate}
  \item เพื่อศึกษาประสิทธิภาพการกำจัดซีโอดีของถังกรองไร้อากาศ
  \item เพื่อหาระยะเวลากักเก็บน้ำที่เหมาะสม
\end{enumerate}

\chapter{การตรวจเอกสาร}

\section{ระบบบำบัดน้ำเสียแบบไร้อากาศ}

\subsection{หลักการทำงาน}

ระบบบำบัดแบบไร้อากาศอาศัยจุลินทรีย์ย่อยสลายสารอินทรีย์ในสภาวะไม่มีออกซิเจน
อัตราการเจริญจำเพาะของจุลินทรีย์อธิบายได้ด้วยสมการของ Monod ดังสมการที่ (\ref{eq:monod})
\begin{equation}
  \mu = \mu_{\max}\,\frac{S}{K_s + S} \label{eq:monod}
\end{equation}
เมื่อ $\mu$ คืออัตราการเจริญจำเพาะ $S$ คือความเข้มข้นของสารอาหาร
และ $K_s$ คือค่าคงที่ครึ่งอิ่มตัว

\subsubsection{ปัจจัยด้านอุณหภูมิ}
อุณหภูมิที่เหมาะสมต่อแบคทีเรียสร้างมีเทนอยู่ในช่วง 30--38 องศาเซลเซียส

\subsubsection{ปัจจัยด้านพีเอช}
ค่าพีเอชที่เหมาะสมอยู่ระหว่าง 6.8--7.2

\subsection{ตัวกลางสำหรับจุลินทรีย์ยึดเกาะ}

ผลการเปรียบเทียบตัวกลางแสดงในตารางที่ \ref{tab:media} และภาพที่ \ref{fig:reactor}

\begin{table}[htbp]
  \caption{คุณสมบัติของตัวกลางชนิดต่าง ๆ ที่ใช้ในถังกรองไร้อากาศ}
  \label{tab:media}
  \centering
  \begin{tabular}{lccc}
    \toprule
    ตัวกลาง & พื้นที่ผิวจำเพาะ & ความพรุน & ราคา \\
            & (m\textsuperscript{2}/m\textsuperscript{3}) & (\%) & (บาท/m\textsuperscript{3}) \\
    \midrule
    เส้นใยไนล่อน & 320 & 96 & 1,200 \\
    พลาสติกริง   & 210 & 92 & 3,500 \\
    หินกรวด      & 60  & 45 & 400 \\
    \bottomrule
  \end{tabular}
  \tablesource{\textcite{udomsak2549}}
\end{table}

\begin{figure}[htbp]
  \centering
  \fbox{\parbox[c][1.6in][c]{3.2in}{\centering (แผนภาพถังกรองไร้อากาศ)}}
  \caption{แผนภาพของถังกรองไร้อากาศที่ใช้เส้นใยไนล่อนเป็นตัวกลาง
           ซึ่งชื่อภาพยาวเกินหนึ่งบรรทัดจะพิมพ์บรรทัดถัดไปตรงกับอักษรตัวแรกของชื่อภาพ}
  \label{fig:reactor}
\end{figure}

\section{งานวิจัยที่เกี่ยวข้อง}

\textcite{chokprasin2003} และ \textcite{nabhadalung2005} ศึกษาเรื่องที่เกี่ยวข้อง
รวมทั้ง \textcite{buddhakulsomsiri2008} และ \textcite{seranevijaikitkhan2008}

\chapter{อุปกรณ์และวิธีการ}

\section{อุปกรณ์}
ถังปฏิกรณ์ทำจากอะคริลิกใส ขนาดเส้นผ่านศูนย์กลาง 15 เซนติเมตร

\section{วิธีการ}
เดินระบบต่อเนื่องเป็นเวลา 180 วัน

\chapter{ผลและวิจารณ์}

\section{ผล}
ระบบเข้าสู่สภาวะคงตัวภายใน 45 วัน

\section{วิจารณ์}
ผลการทดลองสอดคล้องกับ \textcite{peterson2004}

\chapter{สรุปและข้อเสนอแนะ}

\section{สรุป}
ถังกรองไร้อากาศที่ใช้เส้นใยไนล่อนเป็นตัวกลางมีประสิทธิภาพสูงและต้นทุนต่ำ

\section{ข้อเสนอแนะ}
ควรศึกษาการเดินระบบที่อุณหภูมิต่ำเพิ่มเติม

\printbibliography

\appendix
\appendixsection{ผลการวิเคราะห์คุณภาพน้ำ}

\begin{table}[htbp]
  \caption{ค่าซีโอดีของน้ำเข้าและน้ำออกจากระบบ}
  \centering
  \begin{tabular}{ccc}
    \toprule
    วันที่ & น้ำเข้า (mg/L) & น้ำออก (mg/L) \\
    \midrule
    1 & 2,450 & 380 \\
    2 & 2,510 & 342 \\
    \bottomrule
  \end{tabular}
\end{table}

\appendixsection{แบบบันทึกข้อมูล}

แบบบันทึกข้อมูลที่ใช้ในการวิจัย

\begin{vitae}
  \vitaeitem{ชื่อ}{นางสาวซันวานี จิใจ}
  \vitaeitem{เกิดวันที่}{14 กุมภาพันธ์ 2520}
  \vitaeitem{สถานที่เกิด}{อำเภอพญาไท จังหวัดกรุงเทพมหานคร}
  \vitaeitem{ประวัติการศึกษา}{วศ.บ. (วิศวกรรมสิ่งแวดล้อม) มหาวิทยาลัยเกษตรศาสตร์}
  \vitaeitem{ตำแหน่งปัจจุบัน}{วิศวกร}
  \vitaeitem{สถานที่ทำงานปัจจุบัน}{กรมควบคุมมลพิษ กระทรวงทรัพยากรธรรมชาติและสิ่งแวดล้อม}
  \vitaeitem{ผลงานดีเด่นและ/หรือรางวัลทางวิชาการ}{รางวัลวิทยานิพนธ์ดีเด่น พ.ศ. 2549}
  \vitaeitem{ทุนการศึกษาที่ได้รับ}{ทุนผู้ช่วยสอนจากบัณฑิตวิทยาลัย มหาวิทยาลัยเกษตรศาสตร์ (พ.ศ. 2547)}
\end{vitae}

\end{document}
